\documentclass[10pt,a4paper]{amsart}
\usepackage[margin=19mm]{geometry}
\usepackage{amsmath,amssymb,mathtools}
\usepackage[scr=rsfs,cal=euler]{mathalfa}
\usepackage{xfrac}
\usepackage[unicode,hidelinks,bookmarks=false]{hyperref}
\newcommand{\ie}{i.e.\ }
\newcommand{\cf}{cf.\ }
\newcommand{\resp}{respectively\ }
\newcommand{\Subsection}{\subsection}
\providecommand{\nobreakdash}{\nobreak\hbox{-}}
\setlength{\emergencystretch}{2em}
\multlinegap0pt
\usepackage{mathtools}
\usepackage{mathrsfs}
\usepackage{tikz}
\usetikzlibrary{shapes.geometric, arrows.meta, positioning}
\usepackage{orcidlink}
\newtheorem{theorem}{Theorem}[section]
\newtheorem{lemma}[theorem]{Lemma}
\newtheorem{example}[theorem]{Example}
\newtheorem{question}[theorem]{Question}
\newtheorem{corollary}[theorem]{Corollary}
\newtheorem{proposition}[theorem]{Proposition}
\newtheorem{claim}[theorem]{Claim}
\newtheorem{definition}[theorem]{Definition}
\newtheorem{remark}[theorem]{Remark}
\newtheorem{counterexample}[theorem]{Counter-example}
\def\union{\mathop{\cup }}
\def\inter{\mathop{\cap }}
\def\petitoplus{\mathop{\bigoplus}}
\begin{document}
\title[Pego theorem for Hilbert space-valued functions on compact g]{Pego theorem for Hilbert space-valued functions on compact groups}
\author{Anaté Kodjovi Lakmon; Yaogan Mensah}
\date{}
\maketitle
\noindent\emph{Source and licence.} Pego theorem for Hilbert space-valued functions on compact groups. Version arXiv:2608.13142v1 (CC BY 4.0). This material is distributed under the Creative Commons Attribution 4.0 International licence, http://creativecommons.org/licenses/by/4.0/, as recorded in the arXiv OAI licence metadata for this exact version and shown on the abstract page https://arxiv.org/abs/2608.13142v1. Full rights notice: licenses/arxiv-2608.13142-CC-BY-4.0.txt.\par\smallskip
\noindent\textit{Editorial scope.} Counted unit: the original \texttt{theorem} environment \texttt{thm:main-tight} at source lines 1009--1019 together with its complete proof at lines 1022--1094; the delivered document keeps source lines 167--1095 so that every referenced label and every citation resolves inside it. Native editable source is the paper's own concatenated TeX; the AMS wrapper is editorial formatting only. No publisher class, upstream script or unrelated artwork is redistributed.
\begin{theorem}\label{scalar:plancherel}
The Fourier transform is an isometric isomorphism from $L^2(G)$ onto
$\ell^2\text{-}\petitoplus\limits_{\sigma\in\widehat{G}}B_2(H_\sigma)$:
$$
\|f\|_2^2=\sum_{\sigma\in\widehat{G}}d_\sigma\,\|\hat f(\sigma)\|_{B_2(H_\sigma)}^2,$$
where $B_2(H_\sigma)$ is the Hilbert-Schmidt class on $H_\sigma$.
\end{theorem}



For $f,g\in L^1(G)$, the convolution of $f$ by $g$ is defined by 
$$
(f*g)(x):=\int_G f(xy^{-1})\,g(y)\,dm_G(y).$$
The following result (convolution theorem) holds:
$$\widehat{f*g}(\sigma)=\widehat{G}(\sigma)\,
\widehat f(\sigma),\quad\sigma\in\widehat{G}.
$$

Let us also mention that for $y\in G$ and $f\in L^p(G)$, 
$$
\widehat{R_yf}(\sigma)=\sigma(y)\,\widehat f(\sigma),\quad\sigma\in\widehat{G}.
$$


\subsection{Pego theorem on compact groups: the scalar case}

In this subsection, we present Pego's theorem as formulated by Kumar in \cite{Kumar2024}. Prior to that, we introduce the definitions essential for its understanding.

Let $(X, d)$ be a metric space.
\begin{definition}
A subset $A$ of $X$ is said to be  {totally bounded} if, for every $\varepsilon > 0$, there exists a finite set of points $\{x_1, x_2, \dots, x_n\} \subset X$ such that:
$$
A \subset \union\limits_{i=1}^{n} B(x_i, \varepsilon),
$$
where $B(a,r)$ denotes the open ball centered at $a$ with radius $r$.
\end{definition}
\begin{definition}
A subset $A$ of $X$ is said to be {precompact} if its topological closure, denoted as $\overline{A}$, is compact.
\end{definition}

\noindent The Hausdorff criterion for precompactness states that a subset of a metric space is precompact if and only if it is totally bounded. 

\begin{definition}
The set $K\subset L^p(G)$ is said to be uniformly $L^p(G)$-equicontinuous if for every
$\varepsilon>0$ there is an open neighborhood $O$ of $e$ such that $\|R_yf-f\|_{L^p(G)}<\varepsilon$
for $f\in K,\ y\in O$.
\end{definition}


Let $B_p(H_\sigma)$ denotes   the Schatten $p$-class of bounded linear operators on $H_\sigma$ endowed with norm  $\|T\|_{B_p(H_\sigma)}=\left(\text{tr}(|T|^p)\right)^{\frac{1}{p}}$.

\begin{definition}
The set $F\subset\ell^p\text{-}\petitoplus\limits_{\sigma\in \widehat{G}} B_p(H_\sigma)$ is said to have uniform
$\ell^p\text{-}\petitoplus\limits_{\sigma\in \widehat{G}} B_p(H_\sigma)$-decay if for every $\varepsilon>0$ there
is a finite set $S\subset\widehat{G}$ such that  $\|\varphi\|_{\ell^p\text{-}\petitoplus_{\sigma \in \widehat{G}\setminus S}B_p(H_\sigma)}<\varepsilon$ for all $\varphi\in F$. 

\end{definition}

\begin{theorem} \cite[Theorem 1.1]{Kumar2024}\label{thm:kumar}
Let $K$ be a bounded subset of $L^2(G)$. The following assertions are equivalent:
\begin{enumerate}
\item[(i)] $K$
is precompact;
\item[(ii)] $K$ is uniformly $L^2(G)$-equicontinuous;
\item[(iii)] $\widehat K$ has uniform $\ell^2\text{-}\petitoplus\limits_{\sigma \in \widehat{G}} B_2(H_\sigma)$-decay.
\end{enumerate}
\end{theorem}
\section{Pego theorem for Hilbert-space-valued functions}\label{sec: main results}
We now establish several results, leading up to the generalization of  Theorem~\ref{thm:kumar} to Hilbert space-valued functions.


 Let $G$ be a compact group. Fix a  complex Hilbert space $\mathcal{H}$. The Bochner spaces of $p$-strongly integrable $\mathcal{H}$-valued functions on $G$ are denoted by $L^p(G,\mathcal{H}), \, 1\le p<\infty$. The norm on $L^p(G,\mathcal{H})$ is given by 
 
$$ \|f\|_{L^p(G,\mathcal{H})}=\left(\int_G\|f(x)\|_{\mathcal{H}}^pdm_G(x)\right)^{\frac{1}{p}}.$$
In particular, $L^2(G,\mathcal{H})$ is a complex  Hilbert space with the inner product
$$\langle f,g\rangle_{L^2(G,\mathcal{H})}=\int_G\langle f(x),g(x)\rangle_{\mathcal{H}}dm_G(x).$$
\begin{definition}\cite{Assiamoua1989}
For $f\in L^1(G,\mathcal{H})$, the Fourier transform  $\widehat f=(\widehat f(\sigma))_{\sigma\in\widehat{G}}$
is the family of sesquilinear maps $\widehat f(\sigma):H_\sigma\times H_\sigma\to\mathcal{H}$ given by 
\[
\widehat f(\sigma)(\xi,\eta)=\int_G\langle\sigma(x)^*\xi,\eta\rangle_{H_\sigma}\,f(x)\,dm_G(x),
\qquad\xi,\eta\in H_\sigma,
\]
where $\langle \cdot, \cdot\rangle_{H_\sigma}$ is the scalar product in $H_\sigma$.
\end{definition}

For $f\in L^2(G,\mathcal{H})$, the inversion formula holds:
\begin{equation}\label{eq:inversionformula}
f(x)=\sum_{\sigma\in\widehat{G}}d_\sigma\sum_{i=1}^{d_\sigma}\sum_{j=1}^{d_\sigma}u^\sigma_{ij}(x)\widehat f(\sigma)(\xi_j^\sigma,
\xi_i^\sigma),\quad x\in G.
\end{equation}

In the sequel, we will write
$$\varphi(\sigma)_{ij}=\varphi(\sigma)(\xi_j^\sigma,
\xi_i^\sigma).$$ In particular, 
$$\widehat f(\sigma)_{ij}=\widehat f(\sigma)(\xi_j^\sigma,
\xi_i^\sigma).$$

We consider the Assiamoua spaces $\mathscr{S}_p(\widehat{G}, \mathcal{H}),\, 1\le p\le \infty$, defined as follows \cite{Mensah2024}:
\begin{itemize}
\item for $1\le p<\infty$, 
$$\mathscr{S}_p(\widehat{G},\mathcal{H})=\Big\{\varphi=(\varphi(\sigma))_{\sigma\in \widehat{G}}:
\sum_{\sigma\in \widehat{G}} d_\sigma\sum_{i=1}^{d_\sigma}\sum_{j=1}^{d_\sigma}\|\varphi(\sigma)_{ij}\|_\mathcal{H}^p<\infty\Big\},
$$
with the norm 
$$ \|\varphi\|_{\mathscr{S}_p(\widehat{G},\mathcal{H})}=\left(\sum_{\sigma\in \widehat{G}} d_\sigma\sum_{i=1}^{d_\sigma}\sum_{j=1}^{d_\sigma}\|\varphi(\sigma)_{ij}\|_\mathcal{H}^p\right)^{\frac{1}{p}}.
$$
\item $$\mathscr{S}_\infty(\widehat{G},\mathcal{H})=\Big\{  \varphi=(\varphi(\sigma))_{\sigma\in \widehat{G}}: \sup_{\sigma \in \widehat{G}}\|\varphi(\sigma)\| <\infty \Big\},$$

with the norm 

$$\|\varphi\|_{\mathscr{S}_\infty(\widehat{G},\mathcal{H})}=\sup_{\sigma \in \widehat{G}}\|\varphi(\sigma)\|.$$
Here,  
$$\|\varphi(\sigma)\|=\sup\Big\{\|\varphi(\sigma)(\xi,\eta)\|_{\mathcal{H}}: \|\xi\|_{H_\sigma}\leq 1, \|\eta\|_{H_\sigma}\leq 1\Big\}.$$
\end{itemize}

 
\noindent Each $\mathscr{S}_p(\widehat{G},\mathcal{H}), \, 1\leq p\leq \infty$ is a complex Banach
space, and in particular $\mathscr{S}_2(\widehat{G},\mathcal{H})$ is a Hilbert space with the inner product
$$\langle\varphi,\psi\rangle_{\mathscr{S}_2}=\sum_{\sigma\in \widehat{G}} d_\sigma\sum_{i=1}^{d_\sigma}\sum_{j=1}^{d_\sigma}\langle
\varphi(\sigma)_{ij},\psi(\sigma)_{ij}\rangle_\mathcal{H}.$$

\noindent We provide an alternative proof for \cite[Theorem 4.8]{Assiamoua1989} concerning the analogue of the Plancherel theorem for this Fourier transform. The counter-example following  Theorem~\ref{theo:plancherel} illustrates why we restrict our attention to a Hilbert space-valued functions. 
\begin{theorem}\label{theo:plancherel}
Let $G$ be a compact group and let $\mathcal{H}$ be a complex Hilbert space.   The Fourier transformation $\mathscr{F}:L^2(G,\mathcal{H})\to \mathscr{S}_2(\widehat{G},\mathcal{H})$ is an isometric isomorphism.
\end{theorem}


\begin{proof}
Fix an orthonormal basis $\{h_k\}_{k\in \mathfrak{K}}$ of the Hilbert space $\mathcal{H}$. By the Peter-Weyl theorem,  the family $\{\sqrt{d_\sigma}\,u^\sigma_{ij} :
\sigma\in\widehat{G},\ 1\le i,j\le d_\sigma\}$ is an orthonormal basis of $L^2(G)$.  Since
$L^2(G,\mathcal{H})$ is isomorphic to $L^2(G)\otimes\mathcal{H}$ canonically as Hilbert spaces, the family
\[
\big\{\sqrt{d_\sigma}\,u^\sigma_{ij}\otimes h_k \;:\; \sigma\in\widehat{G},\ 1\le i,j\le
d_\sigma,\ k\in \mathfrak{K}\big\}
\]
is an orthonormal basis of $L^2(G,\mathcal{H})$. For $f\in  L^2(G,\mathcal{H})$,  the basis coefficient $c^\sigma_{ij,k}$ of $f$ against the vector $\sqrt{d_\sigma}\,u^\sigma_{ij}\otimes h_k$ is
\begin{equation}\label{Coefficient}
c^\sigma_{ij,k}=\big\langle f,\ \sqrt{d_\sigma}\,u^\sigma_{ij}\otimes h_k
\big\rangle_{L^2(G,\mathcal{H})}
=\sqrt{d_\sigma}\int_G \overline{u^\sigma_{ij}(x)}\,\big\langle f(x),h_k\big\rangle_\mathcal{H}
\,dm_G(x).
\end{equation}

Now, we aim to  relate $c^\sigma_{ij,k}$ to the Fourier coefficient $\widehat{f}(\sigma)_{ij}$.
We have
\[
\big\langle\sigma(x)^*\xi_j^\sigma,\xi_i^\sigma\big\rangle_{H_\sigma}
=\big\langle\xi_j^\sigma,\sigma(x)\xi_i^\sigma\big\rangle_{H_\sigma}
=\overline{\big\langle\sigma(x)\xi_i^\sigma,\xi_j^\sigma\big\rangle}_{H_\sigma}
=\overline{u^\sigma_{ij}(x)}.
\]

Hence, directly from the definition of the Fourier transform, we can write 
\[
\widehat f(\sigma)_{ij}
=\int_G\overline{u^\sigma_{ij}(x)}\,f(x)\,dm_G(x).
\]
Taking the inner product with $h_k$ and compairing with (\ref{Coefficient}), we obtain 
\[
\big\langle\widehat f(\sigma)_{ij},h_k\big\rangle_\mathcal{H}
=\int_G\overline{u^\sigma_{ij}(x)}\,\big\langle f(x),h_k\big\rangle_\mathcal{H}\,dm_G(x)
=\frac{c^\sigma_{ij,k}}{\sqrt{d_\sigma}},
\]
i.e.
\begin{equation}\label{CoefficientC}
c^\sigma_{ij,k}=\sqrt{d_\sigma}\;\big\langle \widehat f(\sigma)_{ij},h_k\big\rangle_\mathcal{H}.
\end{equation}
So the basis coefficients of $f$ are, up to the fixed scalar factor $\sqrt{d_\sigma}$,
exactly the coordinates of the Fourier coefficients $\widehat f(\sigma)_{ij}$ against the basis $\{h_k\}$ of $\mathcal{H}$.

\noindent Applying Parseval's identity for
the orthonormal basis
 $\{\sqrt{d_\sigma}u^\sigma_{ij}\otimes h_k:\sigma\in \widehat{G}, 1\le i,j\le d_\sigma, k\in \mathfrak{K}\}$ of $L^2(G,\mathcal{H})$
by using (\ref{CoefficientC}), we attain
\[
\|f\|_{L^2(G,\mathcal{H})}^2=\sum_{\sigma\in \widehat{G}}\sum_{i=1}^{d_\sigma}\sum_{j=1}^{d_\sigma}\sum_{k\in \mathfrak{K}}|c^\sigma_{ij,k}|^2
=\sum_{\sigma\in \widehat{G}}d_\sigma\sum_{i=1}^{d_\sigma}\sum_{j=1}^{d_\sigma}\sum_{k\in \mathfrak{K}}\big|\langle\widehat f(\sigma)_{ij},h_k\rangle_\mathcal{H}\big|^2.
\]
For each fixed $\sigma,i,j$, Parseval's identity in $\mathcal{H}$ itself, applied to the vector
$\widehat f(\sigma)_{ij}\in\mathcal{H}$ against the orthonormal basis $\{h_k:k\in \mathfrak{K}\}$, gives
\[
\sum_{k\in \mathfrak{K}}\big|\langle \widehat f(\sigma)_{ij},h_k\rangle_\mathcal{H}\big|^2=\|\widehat f(\sigma)_{ij}\|_\mathcal{H}^2.
\]
Substituting back,
\[
\|f\|_{L^2(G,\mathcal{H})}^2=\sum_{\sigma\in \widehat{G}}d_\sigma\sum_{i=1}^{d_\sigma}\sum_{j=1}^{d_\sigma}\|\widehat f(\sigma)_{ij}\|_\mathcal{H}^2
=\|\widehat f\|_{\mathscr{S}_2(\widehat{G},\mathcal{H})}^2.
\]

Therefore, the Fourier transformation $\mathscr{F}$ is an isometry, hence  an  injection. 
\smallskip

Furthermore, given $\varphi\in \mathscr{S}_2(\widehat{G},\mathcal{H})$, that is,  $$\sum_{\sigma\in \widehat{G}}d_\sigma\sum_{i=1}^{d_\sigma}\sum_{j=1}^{d_\sigma}\|\varphi(\sigma)_{ij}\|_\mathcal{H}^2<\infty, $$
 define scalars $q^\sigma_{ij,k}$ by
$q^\sigma_{ij,k}=\sqrt{d_\sigma}\langle\varphi(\sigma)_{ij},h_k\rangle_\mathcal{H}$. By the same computation as above, $$\sum_{\sigma\in \widehat{G}}\sum_{i=1}^{d_\sigma}\sum_{j=1}^{d_\sigma}\sum_{k\in \mathfrak{K}}
|q^\sigma_{ij,k}|^2=\|\varphi\|_{\mathscr{S}_2(\widehat{G},\mathcal{H})}^2<\infty.$$
Since $L^2(G,\mathcal{H})$ is complete,  the series
\[
\sum_{\sigma\in \widehat{G}}\sum_{i=1}^{d_\sigma}\sum_{j=1}^{d_\sigma}\sum_{k\in \mathfrak{K}}q^\sigma_{ij,k}\;\sqrt{d_\sigma}\,u^\sigma_{ij}\otimes h_k
\]
converges in $L^2(G,\mathcal{H})$ to a well-defined element $g$, and  
 $\widehat{G}(\sigma)_{ij}=\varphi(\sigma)_{ij}$ for every $\sigma,i,j$,
i.e.\ $\widehat{G}=\varphi$. Hence, $f\mapsto\widehat f$ is surjective. 
\end{proof}
In Theorem~\ref{theo:plancherel}, the assumption that \(\mathcal{H}\) is a Hilbert space is essential, without which the Fourier transform may fail to be an isometry, as illustrated by the following counter-example.

\begin{counterexample}\label{ex:planch-fails-Cstar}{\rm
Take $G=\mathbb T$ be the 1-dimensional torus (here $\widehat{G}=\mathbb Z$ and the degrees of the characters are equal to 1)  and $\mathcal{H}=\mathbb C^2$ equipped with
the norm $\|(z_1,z_2)\|_{\mathbb C^2}=\max(|z_1|,|z_2|)$. This norm on $\mathbb C^2$ does not satisfy the parallelogram law, so $\left(\mathbb C^2,\|\cdot\|_{\mathbb C^2} \right)$ is not a Hilbert space. 

Let $u_1(x)=e^{ix}$, $u_2(x)=e^{2ix}$ be orthonormal characters of
$\mathbb T$, and set
$$
h(x)=a_1u_1(x)+a_2u_2(x)$$
where $a_1=(1,0)$ and $a_2=(0,1)$.
Then $h(x)=(e^{ix},e^{2ix})$, so $\|h(x)\|_{\mathbb C^2}=\max(1,1)=1$ for every $x\in\mathbb T$,
giving
\[
\|h\|_{L^2(\mathbb T,\mathbb C^2)}^2=1.
\]

\noindent But  the Fourier coefficients computed directly are $\widehat h(u_1)=a_1$, $\widehat h(u_2)=a_2$,
and all other coefficients vanish. So, 
\[
\|\widehat h\|_{\mathscr{S}_2(\mathbb Z,\mathbb C^2)}^2=\|a_1\|_{\mathbb C^2}^2+\|a_2\|_{\mathbb C^2}^2=1+1=2\neq 1
=\|h\|_{L^2(\mathbb T,\mathbb C^2)}^2.
\]
 Therefore, the  Fourier transform is not
an isometry from $L^2(\mathbb T,\mathbb C^2)$ into $\mathscr{S}_2(\mathbb Z,\mathbb C^2)$. 
}\end{counterexample}
Counter-example~\ref{ex:planch-fails-Cstar}  provides a negative answer to Question 2 in \cite{Mensah2024}. 

\smallskip

The following result is the analogue of the Hausdorff-Young inequality. We prove it via an interpolation method, bypassing the theory of Fourier type \cite{GarciaRubio, Peetre}.

\begin{theorem}\label{prop:HY-final}
Let $G$  be a compact group and  $\mathcal{H}$ a complex Hilbert space. Let   $1\le p\le 2$. If $f\in L^p(G,\mathcal{H})$, then
\[
\|\widehat f\|_{\mathscr{S}_{p'}(\widehat{G},\mathcal{H})}\le\|f\|_{L^p(G,\mathcal{H})},\quad\text{where }
\frac1p+\frac1{p'}=1.
\]
\end{theorem}

\begin{proof}
If $f\in L^1(G,\mathcal{H})$, then $\widehat{f}\in \mathscr{S}_\infty(\widehat{G},\mathcal{H})$. Indeed,
for every $\sigma\in\widehat{G}$ and $\xi,\eta\in H_\sigma$ with
$\|\xi\|_{H_\sigma}=\|\eta\|_{H_\sigma}=1$, we have
\begin{align*}
\|\hat f(\sigma)(\xi,\eta)\|_\mathcal{H}&=\Big\|\int_G\langle\sigma(x)^*\xi,\eta\rangle_{H_\sigma}\,
f(x)\,dm_G(x)\Big\|_\mathcal{H}\\
&\le \int_G|\langle\sigma(x)^*\xi,\eta\rangle_{H_\sigma}|
\|f(x)\|_\mathcal{H}dm_G(x)\\
&\le\int_G\|f(x)\|_\mathcal{H}\,dm_G(x)\\
&=\|f\|_{L^1(G,\mathcal{H})},
\end{align*}
using meanwhile the Cauchy-Schwarz inequality and the fact that  $\sigma(x)^*$ is unitary to obtain $|\langle\sigma(x)^*\xi,\eta\rangle_{H_\sigma}|\le1$.
Hence, 
\begin{equation}\label{OneInfinity}
\|\widehat f\|_{\mathscr{S}_\infty(\widehat{G},\mathcal{H})}\le\|f\|_{L^1(G,\mathcal{H})}. 
\end{equation}
Moreover, by Theorem~\ref{theo:plancherel}, we have 
\begin{equation}
\|\widehat f\|_{\mathscr{S}_2(\widehat{G},\mathcal{H})}=\|f\|_{L^2(G,\mathcal{H})}. 
\end{equation}

Now, using  \cite[Theorem 5.1.2]{BerghLofstrom} related to the  interpolation of  Bochner $L^p$-spaces,  we obtain  that if $1\leq p\leq 2$, then  
$$
\|\widehat f\|_{\mathscr{S}_{p'}(\widehat{G},\mathcal{H})}\le\|f\|_{L^p(G,\mathcal{H})},
$$
after noting that  each  space  $\mathscr{S}_p(\widehat{G},\mathcal{H})$ carries a  counting measure 
weighted by $d_\sigma$ on the discrete index set $\{(\sigma,i,j):\sigma\in \widehat{G}, 1\le i,j\le d_\sigma\}$. 
\end{proof}


We now prove the analogue of the inverse Hausdorff-Young inequality.
\begin{theorem}\label{theo:invHY}
Let $G$ be a compact group,  $\mathcal{H}$  a complex Hilbert space  and $1\le p\le2$. The inversion map $\varphi\mapsto\varphi^\vee$, defined by
\begin{equation}
\varphi^\vee(x)=\sum_{\sigma\in\widehat{G}}d_\sigma\sum_{i=1}^{d_\sigma}\sum_{j=1}^{d_\sigma}u^\sigma_{ij}(x)\varphi(\sigma)_{ij},\quad x\in G,
\end{equation}
extends to a bounded linear operator $\mathscr{S}_p(\widehat{G},\mathcal{H})\to L^{p'}(G,\mathcal{H})$, with
\[
\|\varphi^\vee\|_{L^{p'}(G,\mathcal{H})}\le\|\varphi\|_{\mathscr{S}_p(\widehat{G},\mathcal{H})},\qquad
\frac1p+\frac1{p'}=1.
\]
\end{theorem}

\begin{proof}
 Let $\varphi\in\mathscr
S_1(\widehat{G},\mathcal{H})$. Then,  $\sum\limits_{\sigma\in\widehat{G}}d_\sigma\sum\limits_{i=1}^{d_\sigma}\sum\limits_{j=1}^{d_\sigma}\|\varphi(\sigma)_{ij}\|_\mathcal{H}<\infty$.
We have   $|u^\sigma_{ij}(x)|=|\langle\sigma(x)\xi_i^\sigma,
\xi_j^\sigma\rangle_{H_\sigma}|\le\|\sigma(x)\|_{B(H_\sigma)}\|\xi_i^\sigma\|_{H_\sigma}\|\xi_j^\sigma\|_{H_\sigma}=1$
for every $x\in G$, using  the Cauchy-Schwarz inequality, the unitarity of $\sigma(x)$ and the fact that  $\|\xi_j^\sigma\|_\sigma=\|\xi_i^\sigma\|_\sigma
=1$. Hence, for every $x$,
\begin{align*}
\sum_{\sigma\in\widehat{G}}d_\sigma\sum_{i=1}^{d_\sigma}\sum_{j=1}^{d_\sigma}\big\|u^\sigma_{ij}(x)\varphi(\sigma)_{ij}\big\|_\mathcal{H}&\le
\sum_{\sigma\in\widehat{G}}d_\sigma\sum_{i=1}^{d_\sigma}\sum_{j=1}^{d_\sigma}\|\varphi(\sigma)_{ij}\|_\mathcal{H}\\
&=\|\varphi\|_{\mathscr{S}_1(\widehat{G},\mathcal{H})}
<\infty,
\end{align*}
independently of $x$. Since $\mathcal{H}$ is complete, this shows that the series defining
$\varphi^\vee(x)$ converges absolutely in $\mathcal{H}$  and the convergence is
uniform in $x$. A uniform limit of the
partial sums, each being a finite sum of continuous $\mathcal{H}$-valued functions $u^\sigma_{ij}$, is itself continuous. So $\varphi^\vee$ is a well-defined element of $C(G,\mathcal{H})\subset
L^\infty(G,\mathcal{H})$, and
\begin{equation}\label{InfinityOne}
\|\varphi^\vee\|_{L^\infty(G,\mathcal{H})}=\sup_{x\in G}\|\varphi^\vee(x)\|_\mathcal{H}\le
\|\varphi\|_{\mathscr{S}_1(\widehat{G},\mathcal{H})}.
\end{equation}
Furthermore,  by Theorem~\ref{theo:plancherel}  the
map $f\mapsto\widehat f$ is an isometric isomorphism from $L^2(G,\mathcal{H})$ onto $\mathscr{S}_2(\widehat{G},\mathcal{H})$.
Its inverse is exactly the map $\varphi\mapsto\varphi^\vee$ above, restricted to $\mathscr{S}_2(\widehat{G},\mathcal{H})$. Hence
\begin{equation}\label{TwoTwo}
\|\varphi^\vee\|_{L^2(G,\mathcal{H})}=\|\varphi\|_{\mathscr{S}_2(\widehat{G},\mathcal{H})},\quad
\varphi\in\mathscr{S}_2(\widehat{G},\mathcal{H}).
\end{equation}

 By interpolating between (\ref{InfinityOne}) and (\ref{TwoTwo}) based on  \cite[Theorem 5.1.2]{BerghLofstrom},   
we obtain for $1\le p\le 2$,
$$
\|\varphi^\vee\|_{L^{p'}(G,\mathcal{H})}\le\|\varphi\|_{\mathscr{S}_p(\widehat{G},\mathcal{H})},\quad
\varphi\in\mathscr{S}_p(\widehat{G},\mathcal{H}).
$$
\end{proof}


For $K\subset L^p(G,\mathcal{H})$,  set 
$\widehat K=\{\widehat f:f\in K\}$.
Next, we present the two primary definitions of this article.
\begin{definition}\label{def:equicont}
The set $K\subset L^p(G,\mathcal{H})$ is  said to be uniformly $L^p(G,\mathcal{H})$-equicontinuous if for every
$\varepsilon>0$ there is an open neighborhood $O$ of $e$ such that $$\|R_yf-f\|_{L^p(G,\mathcal{H})}
<\varepsilon$$ for all $f\in K,\ y\in O$.
\end{definition}

\begin{definition}\label{def:decay}
The set $F\subset \mathscr{S}_p(\widehat{G},\mathcal{H})$ is said to have uniform $ \mathscr{S}_p(\widehat{G},\mathcal{H})$-decay if for every
$\varepsilon>0$ there is a finite set $S\subset\widehat{G}$ such that  $$\|\varphi\|_{\mathscr{S}_p(\widehat{G}\setminus
S,\mathcal{H})}<\varepsilon$$ for all $\varphi\in F$.
\end{definition}

\bigskip

For $f\in L^p(G,\mathcal{H})$ and for the complex-valued function $k\in L^1(G)$, the convolution $f*k$ is defined by 
\begin{equation}\label{Convolution}
(f*k)(x)=\int_G k(y)f(xy^{-1})
dm_G(y).
\end{equation}
 In Fourier domain, we have, for $\sigma\in\widehat{G}$,
\begin{equation}\label{ConvolutionFourierDomain}
\widehat{f*k}(\sigma)=\widehat k(\sigma)\widehat f(\sigma).
\end{equation}
In coordinates,  Equality~\ref{ConvolutionFourierDomain} turns to $$\widehat{f*k}(\sigma)_{ij}=\sum_{\ell=1}^{d_\sigma}\widehat k(\sigma)_{i\ell}
\,\widehat{f}(\sigma)_{\ell j}.$$
This is an ordinary matrix product of the scalar
matrix $\widehat k(\sigma)$ by the $\mathcal{H}$-valued matrix $\widehat{f}(\sigma)$, with
matrix multiplication acting entrywise via the vector space structure of $\mathcal{H}$. 

\bigskip


\begin{lemma}\label{lem:decay}
Let $G$ be a compact group and $\mathcal{H}$ a complex Hilbert space. Let $K\subset L^2(G,\mathcal{H})$ be uniformly $L^2(G,\mathcal{H})$-equicontinuous. Then $\widehat K$
has uniform $\mathscr{S}_2(\widehat{G},\mathcal{H})$-decay.
\end{lemma}

\begin{proof}
Let $\varepsilon>0$. The uniform $L^2(G,\mathcal{H})$-equicontinuity of $K$ implies the existence of an open neighborhood $U$ of $e$ such that $\|R_yf-f\|_{L^p(G,\mathcal{H})}
<\varepsilon$ for $f\in K,\ y\in U$.  
 Set $O=U\cap U^{-1}$. Then, $O$
is a symmetric open neighborhood  of $e$ 
such that $\|R_yf-f\|_{L^2(G,\mathcal{H})}<\varepsilon$ for $f\in K,\ y\in O$.  Set $k=\displaystyle\frac{1}{m_G(O)}1_O$ where $1_O$ designates the characteristic function of $O$. We have 
 $k\ge0$, $k(y^{-1})=k(y), \, \forall y\in O$, $\displaystyle\int_Gk\,dm_G=1$ and 
$k\in  L^2(G)\subset L^1(G)$ (the inclusion is due to the fact that $m_G$ is a finite measure).

Let $f\in K$. We have  $$(f*k)(x)-f(x)=\int_O\big(R_{y^{-1}}f(x)-f(x)\big)k(y)\,dm_G(y).$$ 
Therefore, 
\begin{align*}
\|f*k-f\|_{L^2(G,\mathcal{H})}&=\left(\int_G\left\|\int_O\big(R_{y^{-1}}f(x)-f(x)\big)k(y)\,dm_G(y)\right\|_{\mathcal{H}}^2dm_G(x)\right)^{\frac{1}{2}}\\
&\le \left[\int_G\left(\int_O k(y)\left\|\big(R_{y^{-1}}f(x)-f(x)\big)\right\|_{\mathcal{H}}dm_G(y)\right)^2dm_G(x)\right]^{\frac{1}{2}}\\
&\leq \int_O k(y) \left[\int_G \left\|\big(R_{y^{-1}}f(x)-f(x)\big)\right\|_{\mathcal{H}}^2 dm_G(x)\right]^{\frac{1}{2}} dm_G(y)\\
&\text{(by the the Minkowski's integral inequality)}\\
&=\int_Ok(y)\,\|R_{y^{-1}}f-f\|_{L^2(G,\mathcal{H})}\,dm_G(y)\\
&<\varepsilon \int_Gk(y)\,dm_G(y)=\varepsilon.
\end{align*}

Then, by Theorem~\ref{theo:plancherel}, 
 applied to $f*k-f$, we have
$$
\|\widehat{f*k}-\widehat f\|_{\mathscr{S}_2(\widehat{G},\mathcal{H})}=\|f*k-f\|_{L^2(G,\mathcal{H})}<\varepsilon,\quad f\in K.
$$
 Fix $\sigma$ in $\widehat{G}$. For each $i,j$, using the triangle inequality in $\mathcal{H}$ and the  Cauchy-Schwarz inequality on the sum over $\ell$, we successively have
 \begin{align*}
\big\|\widehat{f*k}(\sigma)_{ij}\big\|_\mathcal{H}&=\Big\|\sum_{\ell=1}^{d_\sigma}\widehat k(\sigma)_{i\ell}
\widehat{f}(\sigma)_{\ell j}\Big\|_\mathcal{H}\\
&\leq \sum_{\ell=1}^{d_\sigma}|\widehat k(\sigma)_{i\ell}|
\left\|\widehat{f}(\sigma)_{\ell j}\right\|_\mathcal{H}\\
&\le\Big(\sum_{\ell=1}^{d_\sigma}|\widehat k(\sigma)_{i\ell}|^2\Big)^{1/2}
\Big(\sum_{\ell=1}^{d_\sigma}\|\widehat{f}(\sigma)_{\ell j}\|_\mathcal{H}^2\Big)^{1/2}.
\end{align*}
Squaring and  summing over $i,j$  yield:
\begin{equation}\label{Path}
\sum_{i=1}^{d_\sigma}\sum_{j=1}^{d_\sigma}\big\|\widehat{f*k}(\sigma)_{ij}\big\|_\mathcal{H}^2\le\Big(\sum_{i=1}^{d_\sigma}\sum_{j=1}^{d_\sigma}
\|\widehat{f}(\sigma)_{ij}\|_\mathcal{H}^2\Big)\cdot\|\widehat k(\sigma)\|_{B_2(H_\sigma)}^2,
\end{equation}
 
\noindent Furthermore, since $k\in L^2(G)$, the (scalar) Plancherel theorem (Theorem~\ref{scalar:plancherel}) gives
$$\sum_{\sigma\in \widehat{G}}
d_\sigma\|\widehat k(\sigma)\|_{B_2(H_\sigma)}^2=\|k\|_{L^2(G)}^2=1<\infty.$$


Since the series converges, its general term must tend to zero. So,  there is a finite set $S\subset\widehat{G}$ such that $d_\sigma\|\widehat k(\sigma)\|_{B_2(H_\sigma)}^2\le\frac{1}{4}$ for $\sigma\notin S$. This yields $\|\widehat k(\sigma)\|_{B_2(H_\sigma)}^2\le \frac{1}{4}$ since $d_\sigma\ge 1$. 

Multiplying (\ref{Path}) by
$d_\sigma$ and summing over $\sigma\notin S$ yield:
$$
\|\widehat{f*k}\|_{\mathscr{S}_2(\widehat{G}\setminus S,\mathcal{H})}\le\tfrac12\|\widehat f\|_{\mathscr{S}_2(\widehat{G}\setminus S,\mathcal{H})}.$$
Now, 
\begin{align*}
\|\widehat f\|_{\mathscr{S}_2(\widehat{G}\setminus S,\mathcal{H})}&=\|\widehat{f}-\widehat{f*k}+\widehat{f*k}\|_{\mathscr{S}_2(\widehat{G}\setminus S,\mathcal{H})}\\
&\le\|\widehat f-\widehat{f*k}\|_{\mathscr{S}_2(\widehat{G}\setminus S,\mathcal{H})}
+\|\widehat{f*k}\|_{\mathscr{S}_2(\widehat{G}\setminus S,\mathcal{H})}\\
&<\varepsilon +\frac{1}{2}\|\widehat f\|_{\mathscr{S}_2(\widehat{G}
\setminus S,\mathcal{H})}.
\end{align*}
Thus,  $\|\widehat f\|_{\mathscr{S}_2(\widehat{G}\setminus S,\mathcal{H})}<2\varepsilon$ for every $f\in K$, with $S$
independent of $f$. This is the uniform $\mathscr{S}_2(G,\mathcal{H})$-decay of $\widehat K$.
\end{proof}

\begin{lemma}\label{lem:equicont}
Let $G$ be a compact group and $\mathcal{H}$ a complex Hilbert space.  Let $1\le p\le 2$ and  $p'$ be such that $\frac{1}{p}+\frac{1}{p'}=1$.   Let $K\subset L^{p'}(G,\mathcal{H})$ be bounded,
and suppose that $\widehat K$ has uniform $\mathscr{S}_p(\widehat{G},\mathcal{H})$-decay. Then $K$ is uniformly
$L^{p'}(G,\mathcal{H})$-equicontinuous.
\end{lemma}


\begin{proof}
Let $\varepsilon>0$. By uniform $\mathscr{S}_p(\widehat{G},\mathcal{H})$-decay of
$\widehat K$, there is a finite set $S\subset\widehat{G}$ such that
$$
\|\widehat f\|_{\mathscr{S}_p(\widehat{G}\setminus S,\mathcal{H})}<\frac{\varepsilon}{3},\quad f\in K.
$$
Let $P_S$ be the band-limiting projection defined by
$$\widehat{P_Sf}(\sigma)=\left\lbrace\begin{array}{cc}
\widehat f(\sigma), & \text{ if } \sigma\in S\\ 
0, & \text{ otherwise. }
\end{array}\right.$$ 
 Equivalently $f-P_Sf$ has Fourier transform
$\widehat f\cdot 1_{\widehat{G}\setminus S}$ where $1_{\widehat{G}\setminus S}$ is the characteristic function of $\widehat{G}\setminus S$. By Theorem~\ref{theo:invHY} applied to $\varphi=\widehat f\cdot\mathbf1_{\widehat{G}\setminus S}
\in\mathscr{S}_p(\widehat{G},\mathcal{H})$, whose inverse Fourier transform is $f-P_Sf$, we obtain
\begin{equation}\label{RLL}
\|f-P_Sf\|_{L^{p'}(G,\mathcal{H})}\le\|\widehat f\|_{\mathscr{S}_p(\widehat{G}\setminus S,\mathcal{H})}<\frac{\varepsilon}{3},
\quad \forall f\in K. 
\end{equation}

The projection $P_S$ commutes with translation. Indeed, for $y\in G$, $f\in L^{p'}(G,\mathcal{H})$,
$\sigma\in\widehat{G}$, $\xi,\eta\in H_\sigma$,  we have

\begin{align*}
\widehat{R_yf}(\sigma)(\xi,\eta)&=\int_G\langle\sigma(x)^*\xi,\eta\rangle_{H_\sigma}\,
f(xy)\,dm_G(x)\\
&=\int_G\langle\sigma(zy^{-1})^*\xi,\eta\rangle_{H_\sigma}f(z)\,dm_G(z)\\ 
&\text{(by the change of variable }   x=zy^{-1} \\
&\text{ and the invariance of the Haar measure})\\
&=\int_G\langle\sigma(y)\sigma(z)^*\xi,\eta\rangle_{H_\sigma}\,f(z)\,dm_G(z)\\
&=\int_G\langle\sigma(z)^*\xi,\sigma(y)^*\eta\rangle_{H_\sigma}\,f(z)\,dm_G(z)\\
&=\hat f(\sigma)\big(\xi,\sigma(y)^*\eta\big),
\end{align*}
 In particular $\widehat{R_yf}$ is
supported on the same set of $\sigma$'s as $\widehat f$.
Hence,
$$
P_S(R_yf)=R_y(P_Sf),\qquad f\in L^{p'}(G,\mathcal{H}),\ y\in G. 
$$
Now, let us show that $P_Sf$ is continuous, and its finitely many coefficients are uniformly
bounded on $K$. By the inversion formula, restricted to the finite set $S$, we have
$$
P_Sf(x)=\sum_{\sigma\in S}d_\sigma\sum_{i=1}^{d_\sigma}\sum_{j=1}^{d_\sigma}\widehat f(\sigma)_{ij}\,
u^\sigma_{ij}(x),
$$
which is a finite sum of continuous $\mathcal{H}$-valued functions, hence continuous on $G$; by
the fact that a continuous vector-valued function on a compact space is bounded, each term of $P_Sf$, and hence $P_Sf$, is
bounded on $G$. Moreover, by  the H\"older's inequality applied to
$\widehat f(\sigma)_{ij}=\displaystyle\int_G\overline{u^\sigma_{ij}(x)}\,f(x)\,dm_G(x)$ against
$u^\sigma_{ij}\in L^p(G)$ (recall $|u^\sigma_{ij}|\le1$, so $u^\sigma_{ij}\in
L^p(G)$ for every $p$, with $\|u^\sigma_{ij}\|_{L^p(G)}\le1$), we have 
$$
\|\widehat f(\sigma)_{ij}\|_\mathcal{H}\le\|u^\sigma_{ij}\|_{L^p(G)}\,\|f\|_{L^{p'}(G,\mathcal{H})}\le
\|f\|_{L^{p'}(G,\mathcal{H})},\quad\sigma\in S,\ f\in K.
$$
Since $K$ is bounded, say for all $f\in K$, $\|f\|_{L^{p'}(G,\mathcal{H})}\le M$ for some positive constant $M$,  then the coefficients set 
 $\{\widehat f(\sigma)_{ij}:\sigma\in S,\,1\le i,j\le
d_\sigma\}$ is uniformly bounded on $K$ by $M$; that is 

\begin{equation}\label{CoeffBound}
\|\widehat f(\sigma)_{ij}\|_\mathcal{H}\le M,\quad f\in K,\, \sigma\in S,\,1\le i,j\le
d_\sigma.
\end{equation}

Since $G$ is compact and each $u^\sigma_{ij}$ is continuous, $u^\sigma_{ij}$ is
uniformly continuous on $G$ by the Heine-Cantor theorem (a continuous function on a compact topological group, viewed as a compact uniform space, is uniformly continuous). As $S$
is finite, there is a single open neighborhood $O$ of $e$ such that, for every
$y\in O$, $\sigma\in S$, $1\le i,j\le d_\sigma$,
\begin{equation}\label{SupBound}
\sup_{x\in G}\big|u^\sigma_{ij}(xy)-u^\sigma_{ij}(x)\big|<
\frac{\varepsilon}{3M\sum\limits_{\sigma\in S}d_\sigma^3}.
\end{equation}
For $f\in K$ and $y\in O$, we have
\begin{align*}
\big\|R_y(P_Sf)(x)-P_Sf(x)\big\|_\mathcal{H}&=\Big\|\sum_{\sigma\in S}d_\sigma\sum_{i=1}^{d_\sigma}\sum_{j=1}^{d_\sigma}
\widehat f(\sigma)_{ij}\big(u^\sigma_{ij}(xy)-u^\sigma_{ij}(x)\big)\Big\|_\mathcal{H}\\
&\le \sum_{\sigma\in S}d_\sigma\sum_{i=1}^{d_\sigma}\sum_{j=1}^{d_\sigma}\Big\|\widehat f(\sigma)_{ij}\Big\|_\mathcal{H}\big |u^\sigma_{ij}(xy)-u^\sigma_{ij}(x)\big|\\
&<\frac{\varepsilon}{3}, \text{ using (\ref{CoeffBound}) and (\ref{SupBound}).}
\end{align*}
 Integrating  the $p'$-th power of both sides over $G$, we obtain
\begin{equation}\label{RL}
\|R_y(P_Sf)-P_Sf\|_{L^{p'}(G,\mathcal{H})}<\frac{\varepsilon}{3},\quad f\in K,\ y\in O. 
\end{equation}

  We break $R_yf-f$ down into three parts:
$$
R_yf-f=R_y(f-P_Sf)+\big(R_y(P_Sf)-P_Sf\big)-(f-P_Sf).
$$
Thus,  using the fact that right translation is an isometry of $L^{p'}(G,\mathcal{H})$
 (due to the right-invariance of $m_G$) and  the inequalities  (\ref{RLL}) and   (\ref{RL}), we obtain
\begin{align*}
\|R_yf-f\|_{L^{p'}(G,\mathcal{H})}&\le\|R_y(f-P_Sf)\|_{L^{p'}(G,\mathcal{H})}
+\|R_y(P_Sf)-P_Sf\|_{L^{p'}(G,\mathcal{H})}+\|P_Sf-f\|_{L^{p'}(G,\mathcal{H})}\\
&<\frac{\varepsilon}{3}+\frac{\varepsilon}{3}+\frac{\varepsilon}{3}=\varepsilon.
\end{align*}
Hence, $K$ is uniformly
$L^{p'}(G,\mathcal{H})$-equicontinuous.
\end{proof}


\begin{corollary}\label{cor:equiv}
Let $G$ be a compact group and $\mathcal{H}$ a complex Hilbert space. Let $K\subset L^2(G,\mathcal{H})$ be bounded. Then,  $K$ is uniformly $L^2(G,\mathcal{H})$-equicontinuous if and only if
$\widehat K$ has uniform $\mathscr{S}_2(G,\mathcal{H})$-decay.
\end{corollary}

\begin{proof}
 Combining Lemma~\ref{lem:decay}  and  Lemma~\ref{lem:equicont} (restricted  to the case $p=p'=2$)  yields the desired statement.
\end{proof}

\begin{lemma}\label{lemma:translation-continuous}
Let $G$ be a compact group, $\mathcal{H}$ a complex Hilbert space, and $f\in L^2(G,\mathcal{H})$.
Then, the map $G\to L^2(G,\mathcal{H})$, $y\mapsto R_yf$ is continuous. 
\end{lemma}

\begin{proof}
By the construction of the Bochner integral,  continuous $\mathcal{H}$-valued functions are dense
in $L^2(G,\mathcal{H})$. So, given $\delta>0$, choose $g\in C(G,\mathcal{H})$ such that
\begin{equation}\label{DI}
\|f-g\|_{L^2(G,\mathcal{H})}<\delta.
\end{equation}

Since  $R_y$ is an isometry of $L^2(G,\mathcal{H})$ for every $y\in G$,  we have
\begin{equation}\label{DDI}
\|R_yf-R_yg\|_{L^2(G,\mathcal{H})}=\|R_y(f-g)\|_{L^2(G,\mathcal{H})}=\|f-g\|_{L^2(G,\mathcal{H})}<\delta.
\end{equation}



Since $G$ is compact and $g$ is continuous, $g$ is uniformly continuous. Thus, for every
$\eta>0$ there is an open neighborhood $V$ of $e$ such that
$$
\|g(xy)-g(x)\|_\mathcal{H}<\eta,\qquad x\in G,\ y\in V.
$$
 Let  $y\in V$. We have
$$
\|R_yg-g\|_{L^2(G,\mathcal{H})}^2=\int_G\|g(xy)-g(x)\|_\mathcal{H}^2\,dm_G(x)\le\eta^2\,m_G(G)=\eta^2.
$$
Thus, 
\begin{equation}\label{DDDI}
\|R_yg-g\|_{L^2(G,\mathcal{H})}\le\eta,\quad
y\in V.
\end{equation}
Now,  given $\varepsilon>0$, first
apply  (\ref{DI}) with $\delta=\frac{\varepsilon}{3}$ to get $g\in C(G,\mathcal{H})$ such that
$\|f-g\|_{L^2(G,\mathcal{H})}<\frac{\varepsilon}{3}$. Then apply (\ref{DDDI}) with $\eta=\frac{\varepsilon}{3}$ to
get an open neighborhood $V$ of $e$ such that $\|R_yg-g\|_{L^2(G,\mathcal{H})}\le\frac{\varepsilon}{3}$. We now get
\begin{align*}
\|R_yf-f\|_{L^2(G,\mathcal{H})}&\le\|R_yf-R_yg\|_{L^2(G,\mathcal{H})}+\|R_yg-g\|_{L^2(G,\mathcal{H})}+
\|g-f\|_{L^2(G,\mathcal{H})}\\
&<\frac{\varepsilon}{3}+\frac{\varepsilon}{3}+\frac{\varepsilon}{3}=\varepsilon.
\end{align*}
by using (\ref{DDI}) for the first and third terms.
Thus, the map $y\mapsto R_yf$ is continuous at $e$. To check the continuity at any point $y_0\in G$, write
$$R_yf-R_{y_0}f=R_{y_0}\big(R_{y_0^{-1}y}f-f\big)$$
and  use the fact that $\|\cdot\|_{L^2(G,\mathcal{H})}$
is unaffected by the isometry $R_{y_0}$ togheter with the continuity at $e$. The conclusion follows.
\end{proof}

\begin{lemma}\label{lem:PS-finite-dim}
Let $G$ be a compact group, $\mathcal{H}$ a complex Hilbert space, and $S$  a finite subset of $\widehat{G}$.
Let $P_S:L^2(G,\mathcal{H})\to L^2(G,\mathcal{H})$ be the band-limiting projection defined by
$$\widehat{P_Sf}(\sigma)=\left\lbrace\begin{array}{cc}
\widehat f(\sigma), & \text{ if } \sigma\in S\\ 
0, & \text{ otherwise. }
\end{array}\right.$$ 
 Consider the finite number $
N=\sum\limits_{\sigma\in S}d_\sigma^2$.

Then, the range $P_S\big(L^2(G,\mathcal{H})\big)$ is
linearly isometric (up to a fixed rescaling of the norm) to $\mathcal{H}^N$,  the
direct sum of $N$ copies of $\mathcal{H}$.
\end{lemma}

\begin{proof}
 Set $\mathcal{I}=\{(\sigma,i,j): \sigma\in S, 1\le i,j\le d_\sigma\}$. The cardinality of $\mathcal{I}$ is $N$. Enumerate the  triples
$(\sigma,i,j)$ in  $\mathcal{I}$  as $1,\dots,N$. Define
$$
\Psi:\mathcal{H}^N\longrightarrow L^2(G,\mathcal{H}),\quad
\Psi\big((a_{\sigma,i,j})_{(\sigma,i,j)\in \mathcal{I}} \big)=\sum_{\sigma\in S}d_\sigma
\sum_{i=1}^{d_\sigma}\sum_{j=1}^{d_\sigma}u^\sigma_{ij}(\,\cdot\,)a_{\sigma,i,j}.
$$
The map $\Psi$ is manifestly linear, and well-defined into $L^2(G,\mathcal{H})$ since it is a finite sum
of continuous, hence $L^2(G,\mathcal{H})$-valued functions.

\noindent The image of $\Psi$ is exactly $P_S\big(L^2(G,\mathcal{H})\big)$. Indeed, if $g=\Psi(a)$
for some $a\in\mathcal{H}^N$, then
 $\widehat{g}=\widehat{g}\cdot 1_S$; that is  $g=P_Sg$. So,  the image of $\Psi$ lies in
$P_S\big(L^2(G,\mathcal{H})\big)$. Conversely, for any $f\in L^2(G,\mathcal{H})$, setting
$a_{\sigma,i,j}=\hat f(\sigma)_{ij}$ for $\sigma\in S$ gives, by the inversion formula
restricted to $S$,
$$
\Psi(a)=\sum_{\sigma\in S}d_\sigma\sum_{i=1}^{d_\sigma}\sum_{j=1}^{d_\sigma}u^\sigma_{ij}\widehat f(\sigma)_{ij}=P_Sf,
$$
so every element of $P_S\big(L^2(G,\mathcal{H})\big)$ is of the form $\Psi(a)$. Hence,
$\Psi\big(\mathcal{H}^N\big)=P_S\big(L^2(G,\mathcal{H})\big)$.

Let us prove that  $\Psi$ is an  isometry.  By  Theorem~\ref{theo:plancherel} applied to $g=\Psi(a)$, we have 
$$
\|\Psi(a)\|_{L^2(G,\mathcal{H})}^2=\|\widehat{g}\|_{\mathscr{S}_2(G,\mathcal{H})}^2=\sum_{\sigma\in S}
d_\sigma\sum_{i,j}\|a_{\sigma,ij}\|_\mathcal{H}^2.
$$
Define the norm of  $a=(a_{\sigma,i,j})_{(\sigma,i,j)\in \mathcal{I}}\in \mathcal{H}^N$  by:
$$\|a\|_{w}^2=\sum_{\sigma\in S}d_\sigma\sum_{i=1}^{d_\sigma}\sum_{j=1}^{d_\sigma}\|a_{\sigma,ij}\|_\mathcal{H}^2.$$ 

So
\[
\|\Psi(a)\|_{L^2(G,\mathcal{H})}=\|a\|_w,\qquad a\in\mathcal{H}^N,
\]
i.e.\ $\Psi$ is an isometry from $(\mathcal{H}^N,\|\cdot\|_w)$ onto $P_S\big(L^2(G,\mathcal{H})\big)$
with its $L^2(G,\mathcal{H})$-norm. In particular $\Psi$ is injective (isometries are
injective). In conclusion, 
 $P_S\big(L^2(G,\mathcal{H})\big)$ is linearly isometric, up to this fixed rescaling, to
$\mathcal{H}^N$.
\end{proof}

The trivial representation of $G$ is the representation $\sigma_0$ given by  $\sigma_0(x)w=w$ for all $x\in G, \,w\in \mathbb C$. 


\begin{proposition}\label{Claim}
Let $\mathcal{H}$ be a complex Hilbert space,  $a\in \mathcal{H}$, and $f\in L^2(G,\mathcal{H})$ the constant function $f(x)\equiv a$. Then
$\widehat f(\sigma)=0$ for every $\sigma\in\widehat{G}\setminus\{\sigma_0\}$.
\end{proposition}

\begin{proof}
Fix  $\sigma\in\widehat{G}$ with $\sigma\ne\sigma_0$. By definition, for $\xi,\eta\in H_\sigma$,  we have
$$
\widehat f(\sigma)(\xi,\eta)=\int_G\langle\sigma(x)^*\xi,\eta\rangle_{H_\sigma}\,f(x)\,
dm_G(x)=\Big(\int_G\langle\sigma(x)^*\xi,\eta\rangle_{H_\sigma}\,dm_G(x)\Big)\,a.
$$
 It remains to show that scalar integral  $\displaystyle\int_G\langle\sigma(x)^*\xi,
\eta\rangle_{H_\sigma}\,dm_G(x)$ vanishes.

By definition of the scalar Fourier transform
on $G$ applied to the constant function $\mathbf 1(x)\equiv1\in L^1(G)$, we have 
$$
\widehat{\mathbf 1}(\sigma)=\int_G\mathbf 1(x)\,\sigma(x)^*\,dm_G(x)=\int_G\sigma(x)^*\,
dm_G(x),
$$
so $\displaystyle\int_G\langle\sigma(x)^*\xi,\eta\rangle_{H_\sigma}\,dm_G(x)=\langle\widehat{\mathbf1}(\sigma)
\xi,\eta\rangle_{H_\sigma}$, and it suffices to show that $\widehat{\mathbf 1}(\sigma)=0$.

 For any fixed
$y\in G$, we have
\begin{align*}
\sigma(y)\,\hat{\mathbf1}(\sigma)&=\sigma(y)\int_G\sigma(x)^*\,dm_G(x)\\
&=\int_G
\sigma(y)\sigma(x)^*\,dm_G(x)\\
&=\int_G\sigma(y x^{-1})\,dm_G(x)\\
&(\text{by the invariance of } m_G)\\
&=\int_G\sigma(x^{-1})\,dm_G(x)\\
&=\int_G\sigma(x)^*\,dm_G(x)=\widehat{\mathbf 1}(\sigma).
\end{align*}
A similar  computation  gives
$\widehat{\mathbf1}(\sigma)\sigma(y)=\hat{\mathbf1}(\sigma)$. Therefore, 
$\widehat{\mathbf1}(\sigma)$ commutes with $\sigma(y)$ for every $y\in G$. 

 Since $\sigma$ is irreducible, Schur's lemma applied
to the  operator $\widehat{\mathbf1}(\sigma)$  gives
$\widehat{\mathbf1}(\sigma)=c\cdot\mathrm{Id}_{H_\sigma}$ for some scalar $c\in\mathbb C$.
But $\hat{\mathbf 1}(\sigma)=\sigma(y)\hat{\mathbf1}(\sigma)$ for every $y$, thus 
 $c\cdot\mathrm{Id_{H_\sigma}}=c\cdot\sigma(y)$ for every $y\in G$.  If $c\ne0$ this would
give $\sigma(y)=\mathrm{Id}_{H_\sigma}$ for every $y\in G$, i.e.\ $\sigma=\sigma_0$
which contradicts $\sigma\ne\sigma_0$. Hence $c=0$, i.e.
$$
\widehat{\mathbf1}(\sigma)=0,\quad\sigma\in\widehat{G}\setminus\{\sigma_0\}.
$$

\end{proof}


Hereafter is a vector version of the Pego theorem on compact groups (Version 1). 
 
\begin{theorem}\label{thm:main}
Let $G$ be a compact group and  $\mathcal{H}$ a finite-dimensional complex Hilbert space. Let  $K$ be a bounded subset of
$L^2(G,\mathcal{H})$. The following assertions are equivalent: 
\begin{enumerate}
\item[(i)] $K$ is precompact;
\item[(ii)] $K$ is uniformly $L^2(G,\mathcal{H})$-equicontinuous;
\item[(iii)] $\widehat K$ has uniform
$\mathscr{S}_2(\widehat{G},\mathcal{H})$-decay.
\end{enumerate}
\end{theorem}



 
\begin{proof}
By Corollary~\ref{cor:equiv}, assertions (ii) and (iii) are already known to be equivalent. It therefore suffices to prove
(i)$\Rightarrow$(ii) and (iii)$\Rightarrow$(i) to close the cycle and establish the equivalence of all the three assertions.

\medskip
\noindent\textbf{(i) $\Rightarrow$ (ii).}
Let $\varepsilon>0$. Since $K$ is precompact, it is totally bounded, so there exist
$f_1,\dots,f_m\in L^2(G,\mathcal{H})$ such that
$$
K\subset\bigcup_{i=1}^m B\left(f_i,\frac{\varepsilon}{3}\right),
$$
where $B\left(g,r\right)$ is the open ball of centre $g$ and radius $r$.
For $i\in\{1,\dots,m\}$, the map $y\mapsto R_yf_i$ is continuous from $G$ into
$L^2(G,\mathcal{H})$ (Lemma~\ref{lemma:translation-continuous}). Hence,  there is an open neighborhood $O_i$ of $e$ such that
$\|R_yf_i-f_i\|_{L^2(G,\mathcal{H})}<\frac{\varepsilon}{3}$ 
for $y\in O_i$. Let $O= \inter\limits_{i=1}^mO_i$,
an open neighborhood of $e$.

Let $f\in K$ and $y\in O$. Choose $i$ such that $\|f-f_i\|_{L^2(G,\mathcal{H})}<\frac{\varepsilon}{3}$.
Since $R_y$ is an isometry of $L^2(G,\mathcal{H})$,
$\|R_yf-R_yf_i\|_{L^2(G,\mathcal{H})}=\|f-f_i\|_{L^2(G,\mathcal{H})}<\frac{\varepsilon}{3}$. So, by the triangle
inequality
$$
\|R_yf-f\|_{L^2(G,\mathcal{H})}\le\|R_yf-R_yf_i\|_{L^2(G,\mathcal{H})}+\|R_yf_i-f_i\|_{L^2(G,\mathcal{H})}+
\|f_i-f\|_{L^2(G,\mathcal{H})}<\varepsilon.
$$
Therefore,
$K$ is uniformly $L^2(G,\mathcal{H})$-equicontinuous.

\medskip
\noindent\textbf{(iii) $\Rightarrow$ (i).}
Let $\varepsilon>0$. By (iii), choose a finite set $S\subset\widehat{G}$ with
$\|\widehat f\|_{\mathscr{S}_2(\widehat{G}\setminus S,\mathcal{H})}<\frac{\varepsilon}{2}$ for all $f\in K$. Let
$P_S$ be the band-limiting projection defined in  Lemma~\ref{lem:PS-finite-dim}.  By Theorem~\ref{theo:plancherel}, we have
$$
\|f-P_Sf\|_{L^2(G,\mathcal{H})}=\|\widehat f-\widehat{P_Sf}\|_{\mathscr{S}_2(\widehat{G},\mathcal{H})}=
\|\widehat f\|_{\mathscr{S}_2(\widehat{G}\setminus S,\mathcal{H})}<\frac{\varepsilon}{2},\quad f\in K.
$$

\noindent By Lemma~\ref{lem:PS-finite-dim}, the
range $P_S\big(L^2(G,\mathcal{H})\big)$ is isomorphic to a finite direct sum of copies of $\mathcal{H}$.  Since $\dim\mathcal{H}<\infty$ by hypothesis, this finite direct sum is finite-dimensional. Moreover, the set $\{P_Sf:f\in K\}$ is a bounded subset of this finite-dimensional space because $P_S$
is an orthogonal projection on $L^2(G,\mathcal{H})$, so $\|P_Sf\|_{L^2(G,\mathcal{H})}\le
\|f\|_{L^2(G,\mathcal{H})}$, and $K$ is bounded by hypothesis. A bounded subset of a
finite-dimensional normed space is precompact according to the Bolzano-Weierstrass theorem.  Thus, 
$\{P_Sf:f\in K\}$ is totally bounded: there exists $\{g_1,\dots,g_k\}\subset
P_S\big(L^2(G,\mathcal{H})\big)$ such that 
$$\{P_Sf:f\in K\}\subset \union\limits_{j=1}^k B\left(g_j,\frac{\varepsilon}{2}\right).$$


For each $f\in K$, there exists  $g_j\in \{g_1,\dots,g_k\}$ such that $\|P_Sf-g_j\|_{L^2(G,\mathcal{H})}<\frac{\varepsilon}{2}$. By the
triangle inequality, 
\begin{align*}
\|f-g_j\|_{L^2(G,\mathcal{H})}&\le\|f-P_Sf\|_{L^2(G,\mathcal{H})}+\|P_Sf-g_j\|_{L^2(G,\mathcal{H})}\\
&<\frac{\varepsilon}{2}+\frac{\varepsilon}{2}=\varepsilon.
\end{align*}
Thus $K\subset\union\limits_{j=1}^kB(g_j,\varepsilon)$; that is,  $K$ is totally bounded in
$L^2(G,\mathcal{H})$. Hence, $K$ is precompact.
%\medskip
\end{proof}

\begin{remark}{\rm
When $\dim\mathcal{H}=1$, Theorem~\ref{thm:main} reduces to the Pego theorem proved by  Kumar (Theorem~\ref{thm:kumar}).
The finite-dimensionality of $\mathcal{H}$ is used only in the last paragraph of
(iii)$\Rightarrow$(i); it cannot be dropped without an additional hypothesis, since a
bounded set in an infinite-dimensional Hilbert space need not be precompact.
}\end{remark}



\begin{counterexample}\label{ex:counterexample}{\rm
Let $G$ be a compact group and  let $\mathcal{H}$ be an {\it infinite-dimensional} Hilbert space, with orthonormal sequence
$\{e_n\}_{n\ge1}$, and let $f_n\in L^2(G,\mathcal{H})$ be the constant function
$f_n(x)\equiv e_n,\quad x\in G$. Set $K=\{f_n : n\ge1\}$.

\begin{itemize}
\item $K$ is bounded: $\|f_n\|_{L^2(G,\mathcal{H})}=\|e_n\|_\mathcal{H}=1$ for every $n$.

\item $K$ is uniformly $L^2(G,\mathcal{H})$-equicontinuous: since $f_n$ is
constant, $R_yf_n=f_n$ for every $y\in G$. So,  Definition~\ref{def:equicont} holds with
$O=G$ for every $\varepsilon>0$.

\item $\widehat K$ has uniform $\mathscr{S}_2(\widehat{G},\mathcal{H})$-decay: a constant
function has Fourier transform supported only on the trivial representation
$\sigma_0\in\widehat{G}$ (see Proposition~\ref{Claim}), so $\widehat f_n(\sigma)=0$ for all $\sigma\ne\sigma_0$ and all $n$. Thus, 
$$
\|\widehat f_n\|_{\mathscr{S}_2(\widehat{G}\setminus\{\sigma_0\},\,\mathcal{H})}=0,\quad n\ge1,
$$
and Definition~\ref{def:decay} holds with the single finite set $S=\{\sigma_0\}$.

\item $K$ is not precompact: for $n\ne m$,
$$
\|f_n-f_m\|_{L^2(G,\mathcal{H})}=\|e_n-e_m\|_\mathcal{H}=\sqrt2.
$$
 Hence, $K$ has no Cauchy, and so no convergent, subsequence.
\end{itemize}
\noindent  Thus conditions (ii) and (iii) of Theorem~\ref{thm:main} both hold for $K$ while
condition (i) fails: the equivalence breaks down as soon as $\dim\mathcal{H}=\infty$.
}\end{counterexample}

\begin{remark}{\rm 
  When $\dim\mathcal{H}=\infty$, a bounded
set of values need not be precompact, and no amount of control in the $G$-direction
can repair that. Therefore an additional hypothesis on the  set of values in $\mathcal{H}$ would be
needed to restore the equivalence. To this end, we introduce the uniform thinness hypothesis to obtain an infinite-dimensional version of the Pego theorem. 
}\end{remark}

\begin{definition}\label{def:tight}
$K\subset L^2(G,\mathcal{H})$ is said to be uniformly tight if for every $\varepsilon>0$ there is a
finite-dimensional subspace $V\subset\mathcal{H}$ such that
$$
\|f - Q_V\circ f\|_{L^2(G,\mathcal{H})}<\varepsilon,\quad f\in K,
$$
where $Q_V:\mathcal{H}\to V$ is the orthogonal projection on $V$.
\end{definition}



Hereafter is a vector version of the Pego theorem on compact groups (Version 2).


\begin{theorem}\label{thm:main-tight}
Let $\mathcal{H}$ be a complex Hilbert space (not assumed finite-dimensional), and let
$K\subset L^2(G,\mathcal{H})$ be bounded and uniformly tight. Then, the following assertions are
equivalent:
\begin{enumerate}
\item[(i)] $K$ is precompact;
\item[(ii)] $K$ is uniformly
$L^2(G,\mathcal{H})$-equicontinuous;
\item[(iii)] $\widehat K$ has uniform $\mathscr{S}_2(\widehat{G},\mathcal{H})$-decay.
\end{enumerate} 
\end{theorem}

\begin{proof}
By Corollary~\ref{cor:equiv}, assertions (ii) and (iii) are already known to be equivalent. It
therefore suffices to prove (i)$\Rightarrow$(ii) and (iii)$\Rightarrow$(i).

\medskip
\noindent\textbf{(i) $\Rightarrow$ (ii).}
This implication does not use tightness or any dimension hypothesis on $\mathcal{H}$. Let $\varepsilon>0$. Since $K$ is precompact, it is totally bounded; that is,  there exist
$f_1,\dots,f_m\in L^2(G,\mathcal{H})$ with $K\subset\union\limits_{i=1}^mB(f_i,\frac{\varepsilon}{3})$. By
Lemma~\ref{lemma:translation-continuous}, for each $i$ the map $y\mapsto R_yf_i$
is continuous at $e$, so there is an open neighborhood $O_i$ of $e$ such that
$\|R_yf_i-f_i\|_{L^2(G,\mathcal{H})}<\frac{\varepsilon}{3}$ for $y\in O_i$. Let $O=\inter\limits_{i=1}^mO_i$.
For $f\in K$ and $y\in O$, there exists $i\in\{1,\cdots,m\}$ with $\|f-f_i\|_{L^2(G,\mathcal{H})}<\frac{\varepsilon}{3}$.
By the triangle inequality and the fact  that $R_y$ is an isometry of $L^2(G,\mathcal{H})$, we have
\begin{align*}
\|R_yf-f\|_{L^2(G,\mathcal{H})}&\le\|R_yf-R_yf_i\|_{L^2(G,\mathcal{H})}+\|R_yf_i-f_i\|_{L^2(G,\mathcal{H})}+
\|f_i-f\|_{L^2(G,\mathcal{H})}\\
&<\frac{\varepsilon}{3}+\frac{\varepsilon}{3}+\frac{\varepsilon}{3}=\varepsilon.
\end{align*}
Hence, $K$ is uniformly $L^2(G,\mathcal{H})$-equicontinuous.

\medskip
\noindent\textbf{(iii) $\Rightarrow$ (i).}
This is where boundedness alone is not enough when $\dim\mathcal{H}=\infty$
(Counter-example~\ref{ex:counterexample}), and where uniform tightness supplies exactly the missing control in $\mathcal{H}$.

Let $\varepsilon>0$.
 By hypothesis, there exists a  finite set 
$S\subset\widehat{G}$ with $\|\widehat f\|_{\mathscr{S}_2(\widehat{G}\setminus S,\mathcal{H})}<\frac{\varepsilon}{4}$ for
all $f\in K$. Let $P_S$ be the band-limiting projection. By Theorem~\ref{theo:plancherel}, we have 
$$
\|f-P_Sf\|_{L^2(G,\mathcal{H})}=\|\widehat f\|_{\mathscr{S}_2(\widehat{G}\setminus S,\mathcal{H})}<
\frac{\varepsilon}{4},\quad f\in K.
$$
By uniform tightness of
$K$, there exists a finite-dimensional subspace $V\subset\mathcal{H}$ such that
$$
\|f-Q_V\!\circ f\|_{L^2(G,\mathcal{H})}<\frac{\varepsilon}{4},\quad f\in K,
$$
where $Q_V:\mathcal{H}\to V$ is the orthogonal projection on $V$. Since $Q_V$ is a bounded linear map on
$\mathcal{H}$, it commutes with the Bochner integral defining $\widehat f(\sigma)$. Indeed,
\begin{align*}
\widehat{Q_V\circ f}(\sigma)(\xi,\eta)&=\int_G\langle \sigma(x)^* \xi,\eta\rangle_{H_\sigma}Q_V(f(x))dx\\
&=Q_V\left(\int_G\langle \sigma(x)^* \xi,\eta\rangle_{H_\sigma}f(x)dx\right)\\
&=Q_V(\widehat f(\sigma)(\xi,\eta)).
\end{align*}
Hence, $\widehat{Q_V\circ f}(\sigma)=Q_V\big(\widehat f(\sigma)\big)$. Applying the definition of $P_S$, we obtain
$P_S(Q_V\!\circ f)=Q_V\!\circ(P_Sf)$. Since $Q_V$ is an orthogonal projection
(hence norm-nonincreasing), we have 
\begin{align*}
\|Q_V\!\circ f-P_S(Q_V\!\circ f)\|_{L^2(G,\mathcal{H})}&=\|\widehat{Q_V\circ f}\|_{\mathscr{S}_2(\widehat{G}\setminus
S,\mathcal{H})}\\
&\le\|\widehat f\|_{\mathscr{S}_2(\widehat{G}\setminus S,\mathcal{H})}<\frac{\varepsilon}{4},\quad f\in K.
\end{align*}

For $f\in K$, we have 
\begin{align*}
\|f-Q_V(P_Sf)\|_{L^2(G,\mathcal{H})}&\le\|f-Q_V\!\circ f\|_{L^2(G,\mathcal{H})}+\|Q_V\!\circ f-Q_V(P_Sf)\|_{L^2(G,\mathcal{H})}\\
&<\frac{\varepsilon}{4}+\frac{\varepsilon}{4}=\frac{\varepsilon}{2},\quad f\in K.
\end{align*}

\noindent The set $\mathcal{A}:=\{Q_V(P_Sf):f\in K\}$ lies in the linear span of the
finitely many functions $u^\sigma_{ij}\otimes v$ with $\sigma\in S, \, 1\leq i,j\leq d_\sigma$ and $v$ ranging over a basis of the finite-dimensional subspace $V$. Therefore,  $\mathcal{A}$ embeds in a finite-dimensional subspace of $L^2(G,\mathcal{H})$. Since $P_S,Q_V$ are both norm-nonincreasing and $K$ is bounded,
$\mathcal{A}$ is a bounded subset of this finite-dimensional space, hence precompact,  i.e.\  totally bounded. There exist $\{g_1,\dots,g_k\}\subset L^2(G,\mathcal{H})$ such that $\mathcal{A}\subset \union\limits_{i=1}^kB(g_i,\frac{\varepsilon}{2})$. 

 For $f\in K$, choose $g_j$ with $\|Q_V(P_Sf)-g_j\|_{L^2(G,\mathcal{H})}
<\frac{\varepsilon}{2}$.  Then, 
\begin{align*}
\|f-g_j\|_{L^2(G,\mathcal{H})}&\le\|f-Q_V(P_Sf)\|_{L^2(G,\mathcal{H})}+
\|Q_V(P_Sf)-g_j\|_{L^2(G,\mathcal{H})}\\
&<\frac{\varepsilon}{2}+\frac{\varepsilon}{2}=\varepsilon.
\end{align*}
 Hence, $K$ is totally bounded, i.e.\ precompact.
\end{proof}


\begin{thebibliography}{99}
\bibitem[Assiamoua1989]{Assiamoua1989} V.S.K. Assiamoua, A. Olubummo, Fourier-Stieltjes transform of vector-valued measures on compact groups, Acta Sci. Math. (Szeged), {\bf 53} (1989), 301-307.
\bibitem[BerghLofstrom]{BerghLofstrom} J. Bergh and J. L\"ofstr\"om, Interpolation Spaces: An Introduction, Grundlehren der Mathematischen Wissenschaften \textbf{223}, Springer-Verlag, Berlin-New York, 1976.
\bibitem[GarciaRubio]{GarciaRubio} J. Garc\'ia-Cuerva and J.L. Rubio de Francia, Weighted Norm Inequalities and Related Topics, North-Holland Mathematics Studies \textbf{116}, North-Holland, Amsterdam, 1985.
\bibitem[Kumar2024]{Kumar2024} M. Kumar, \emph{Pego theorem on compact groups}, Pacific J. Math. \textbf{328}(1) (2024), 137-143.
\bibitem[Mensah2024]{Mensah2024} Y. Mensah, Vector Fourier analysis on compact groups and Assiamoua spaces, Methods Funct. Anal. Topol., {\bf 30}(3-4) (2024), 147-154.
\bibitem[Peetre]{Peetre} J. Peetre, Sur la transformation de Fourier des fonctions \`a valeurs vectorielles, Rend. Sem. Mat. Univ. Padova \textbf{42} (1969), 15-26.
\end{thebibliography}
\end{document}
